\section{The question}\label{sec:intro}

Coding agents spend money in one place: calls to a large \term{host model}. Every turn re-sends the conversation, and
a session of a dozen turns can cost several dollars. Between those calls sit many small choices. Should this session
run on the expensive model or a cheaper one? May the agent read a file without asking the model first? How much
should the model ``think''? A \term{decision model} (``system~1'') answers such questions: it takes a short state
and a typed question and returns a probability for each allowed answer, in tens to hundreds of milliseconds. Jev,
by Typesafe, is one; GPT-6 Luna and Sol, Cloudflare's Clef and Clef-Flash, and small local models (Qwen3, Laya) can
play the same role.

Our question was simple, and we followed it where it went: \emph{what does a decision model buy an agent harness in
quality, cost and speed?}

\begin{keyidea}
A decision model is only worth what its decisions change. It pays off when a cheap, typed decision changes an
expensive action, and it adds nothing when a fixed rule makes the same choice. Most of this report is about telling
those two cases apart with measurements rather than estimates.
\end{keyidea}

\paragraph{How the story goes.} We first looked at a week of live telemetry (\cref{sec:observatory}). We then built
tools to grade decision models against labelled cases and real traces (\cref{sec:judges}), raised the difficulty to
full paired agent sessions (\cref{sec:paired}), and reached a plateau where the question became how to configure the
model rather than which one to pick (\cref{sec:plateau}). \Cref{sec:recommend} says what to set, per host model and
per harness; \cref{sec:limits} says what we did not finish.

\begin{figure}[htbp]
\centering
% The study program as a timeline: stage, size, spend (all from macros).
\begin{tikzpicture}[
  stage/.style={draw=okblue, fill=okblue!6, rounded corners=2pt, align=center, font=\footnotesize,
    text width=2.8cm, minimum height=1.55cm, inner sep=3pt},
  arr/.style={-{Stealth[length=5pt]}, okgray, line width=1pt}]
\node[stage] (a) {\textbf{1. Observatory}\\ \ObEvents\ events\\ no model calls};
\node[stage, right=0.42cm of a] (b) {\textbf{2. Judge benchmark}\\ \DevN\,+\,\HoldN\ cases\\ \$\LoggedSpendTotal};
\node[stage, right=0.42cm of b] (c) {\textbf{3. Real decisions}\\ \TrN\ traced cases\\ \$\TrSpendTotal; Clef \$\ClSpend};
\node[stage, right=0.42cm of c] (d) {\textbf{4. Caching survey}\\ \CaPairs\ matched pairs\\ no model calls};
\node[stage, below=0.9cm of d] (e) {\textbf{5. Paired campaign}\\ \NSessions\ sessions\\ \$\SpendLedger};
\node[stage, left=0.42cm of e] (f) {\textbf{6. Effort follow-ups}\\ \EcPairs\,+\,\EfPairs\ pairs\\ \$\EcSpend\,+\,\$\EfSpend};
\node[stage, left=0.42cm of f] (g) {\textbf{7. Counterfactual}\\ priced policies\\ no model calls};
\node[stage, left=0.42cm of g, draw=okorange, fill=okorange!8] (h) {\textbf{8. S1 holdout}\\ \SOneSessions\ sessions\\ \$\SSpend};
\draw[arr] (a) -- (b); \draw[arr] (b) -- (c); \draw[arr] (c) -- (d); \draw[arr] (d) -- (e);
\draw[arr] (e) -- (f); \draw[arr] (f) -- (g); \draw[arr] (g) -- (h);
\end{tikzpicture}

\caption{The study program in order: what each stage measured, how big it was and what it cost in model calls.}
\label{fig:timeline}
\howtoread{Follow the arrows. Each stage answered a question the previous one raised: telemetry could not say whether
decisions were right (so we built benchmarks), benchmarks could not say what a decision saves (so we ran whole
sessions), and whole sessions raised the effort and decider questions that the follow-ups and S1 settled.}
\end{figure}

\paragraph{How to read the evidence labels.} Results marked \prereg\ test hypotheses fixed before the data;
\screen\ marks exploratory results. Cost ratios are \term{geometric means} of per-pair ratios unless labelled
otherwise, and a 95\,\% interval is given wherever one exists. Every number is generated from committed evidence
files. The full detail of each study, with its own figures and its peer review, is in the companion technical
report v2 (\path{docs/papers/2026-10-02-paired-measurement/}), cited here as ``the companion report''.


\begin{table}[htbp]
\centering\small
\caption{Findings at a glance. ``Preregistered'' means the hypothesis and its threshold were fixed before the data;
``replicated'' means an independent set of scenarios gave the same answer; ``exploratory'' findings are hypotheses.}
\label{tab:findings}
\begin{tabular*}{\textwidth}{@{\extracolsep{\fill}}>{\raggedright\arraybackslash}p{0.31\textwidth} >{\raggedright\arraybackslash}p{0.26\textwidth} l >{\raggedright\arraybackslash}p{0.16\textwidth}@{}}
\toprule
Finding & Number (95\,\% CI unless noted) & Study & Evidence \\
\midrule
Decision calls are fast next to host calls & Jev p50 \ObJevPfifty{}\,ms; host p50 \ObHostPfifty{}\,s & observatory & measured \\
Shadow ``agreement'' was mostly unmatchable & \ObUnmatchN{} of \ObAllN{} could not match & observatory & measured \\
Jev accurate within noise of the best, cheapest & \JevHoldAccK{}/\HoldN{} correct & judge benchmark & preregistered \\
No judge useful on real read decisions & \TrNUseful{} of \NumBase{} judges & trace study & preregistered \\
Routing on Fable saves & \CfFableStickyPairTwo{}\X\ (\CfFableStickyPairCITwo{}) & main-v1 & preregistered \\
Routing on Fable saves, fresh scenarios & 0.581\X\ (0.542--0.625) & S1, H1 & preregistered, replicated \\
Routing on Opus costs more & 1.255\X\ (1.197--1.321) & S1, H5 & preregistered, replicated \\
Rule R* equals the decision model at session start & 1.008\X\ (90\,\% 1.000--1.019) & S1, H2 & preregistered \\
No bundle overhead on Opus & 0.999\X\ (90\,\% 0.984--1.017) & S1, H3 & preregistered \\
Medium effort: Sonnet & \EcRatio{}\X\ (\EcCI{}) & follow-up 1 & preregistered, provenance-limited \\
Medium effort: Fable & \EfRatio{}\X\ (\EfCI{}) & follow-up 2 & preregistered \\
Medium effort: Opus, no saving & 0.992\X\ (0.974--1.009) & S1, H4 & preregistered \\
Review/explain lose quality on Sonnet & lower bound $-$0.063 & S1, H7 & preregistered \\
Live decide-once equals its prediction & 1.014\X\ (90\,\% 0.996--1.032) & S1, H6 & preregistered \\
Effort switches rewrite the cache & \AoEcX{}\X\ cache writes & main-v1 re-analysis & exploratory \\
Keyword proxy keeps quality at a cost & \KpProxyRatio{}\X, turn-pass \KpProxyDtp{} & S1 replay & exploratory \\
\bottomrule
\end{tabular*}

\end{table}
